% Zdroj: PDF priklady-leto-2024.pdf, fyzická str. 25. Značka: specializace UI-SU. Zařazeno: S3.
\section{Metoda nejmenších čtverců}
\szzotazka{léto 2024}{31}{Metoda nejmenších čtverců}

\noindent\textbf{Zadání.}\par
Mějme data s jedním číselným atributem $x_1$ a číselnou cílovou hodnotou $y$ zadaná v tabulce níže.
Pro tato data chceme najít lineární model metodou nejmenších čtverců.

\begin{center}
\begin{tabular}{c|c}
\toprule
$x_1$ & $y$ \\
\midrule
1 & 1 \\
2 & 3 \\
3 & 2 \\
\bottomrule
\end{tabular}
\end{center}

\begin{enumerate}
  \item Popište lineární model a metodu nejmenších čtverců. Jakým způsobem se touto metodou
    odhadují koeficienty modelu?
  \item Odhadněte koeficienty modelu pro data zadaná v tabulce výše.
  \item Spočítejte střední kvadratickou chybu získaného modelu.
\end{enumerate}

\medskip
\noindent\textbf{Řešení.} \textit{(V~PDF bez náčrtu řešení; vlastní vzorové řešení, výpočet ověřen numericky.)}\par
\begin{enumerate}[nosep]
  \item Lineární model $\hat y=\beta_0+\beta_1 x$; metoda nejmenších čtverců volí koeficienty tak, aby \emph{součet čtverců reziduí} $\sum_i(\hat y_i-y_i)^2$ byl nejmenší (řeší se normální rovnice $\beta=(X^\top X)^{-1}X^\top y$, pro jednu proměnnou stačí uzavřené vzorce).
  \item S~průměry $\bar x=2$, $\bar y=2$: $\beta_1=\dfrac{\sum_i(x_i-\bar x)(y_i-\bar y)}{\sum_i(x_i-\bar x)^2}=\dfrac{1}{2}$, $\beta_0=\bar y-\beta_1\bar x=1$. Tedy $\hat y=1+\tfrac12 x$.
  \item Predikce $(1{,}5;\,2;\,2{,}5)$, rezidua $(-0{,}5;\,1;\,-0{,}5)$, $\mathrm{MSE}=\tfrac13\big(0{,}25+1+0{,}25\big)=0{,}5$.
\end{enumerate}
